\chapter{Defense in Depth: Seccomp}

So far, we have used Falco exclusively for \textbf{detection} -- it observes attacks and generates alerts, but it does not prevent them. Runtime detection is a critical capability, but a mature security posture requires \textbf{prevention} as well. In this chapter, we introduce \textbf{Seccomp} (Secure Computing Mode), a Linux kernel mechanism that restricts the set of system calls a container is allowed to make.

The combination of Falco and Seccomp embodies the \textit{defense in depth} principle:
\begin{itemize}
    \item \textbf{Seccomp} operates at the kernel level and \textit{prevents} specific dangerous syscalls from being executed at all -- the attack cannot succeed.
    \item \textbf{Falco} operates at the kernel level and \textit{detects} anomalous behaviors -- even when an attack is blocked, the attempt itself can still be logged.
\end{itemize}

\section{Seccomp in Kubernetes}

A Seccomp profile is a JSON file that defines an \textit{allowlist} or \textit{denylist} of Linux system calls. When applied to a pod, the Linux kernel enforces this profile on all processes in the container. If a process attempts a forbidden syscall, the kernel returns an error (typically \texttt{EPERM}) instead of executing it.

Seccomp profiles are applied via the \texttt{securityContext.seccompProfile} field in the pod specification. In a kubeadm cluster, custom profiles must be present on the worker node at the path:

\begin{center}
\texttt{/var/lib/kubelet/seccomp/<profile-name>.json}
\end{center}

Navigate to the seccomp profiles directory and copy the profile to the worker node using \texttt{scp} from the control plane, renaming it to the generic name \texttt{seccomp-profile.json} expected by our deployment manifest:

\begin{lstlisting}[language=bash]
cd material/seccomp_profiles
# Run on the control plane node:
scp seccomp-profile-connect.json user@<worker-node-ip>:/tmp/
# Then SSH to the worker and move it into place as seccomp-profile.json:
ssh user@<worker-node-ip> sudo mv /tmp/seccomp-profile-connect.json /var/lib/kubelet/seccomp/seccomp-profile.json
\end{lstlisting}

\begin{tcolorbox}[colback=blue!5!white,colframe=blue!75!black,title=Info]
The Seccomp profile files (\texttt{seccomp-profile-connect.json} and \texttt{seccomp-profile-ptrace.json}) are included in the \texttt{material/seccomp\_profiles} directory.
\end{tcolorbox}

\section{Blocking the Reverse Shell}

The reverse shell requires the containerized process to make an outbound TCP connection to the attacker. This connection is established via the \texttt{connect()} system call. By denying this single syscall, we can prevent the reverse shell from succeeding entirely.

The profile \texttt{seccomp-profile-connect.json} has the following structure:

\begin{verbatim}
{
  "defaultAction": "SCMP_ACT_ALLOW",
  "syscalls": [
    { "names": ["connect"], "action": "SCMP_ACT_ERRNO" }
  ]
}
\end{verbatim}

This profile allows all syscalls by default, but explicitly returns \texttt{ERRNO} for any call to \texttt{connect}. 

The instructor has already provided a pre-configured seccomp-enabled deployment manifest under the name \texttt{zipapp-deployment-seccomp.yaml} in the \texttt{manifests} folder. This manifest is already configured to mount the custom Seccomp profile named \texttt{seccomp-profile.json} at the container level.

Apply the manifest and wait for the rollout:

\begin{lstlisting}[language=bash]
kubectl apply -f ../manifests/zipapp-deployment-seccomp.yaml
kubectl rollout status deployment/zipapp
\end{lstlisting}

Now attempt the reverse shell attack again (upload \texttt{reverse\_shell.zip} and click \textbf{``Download All''}). The web application will return an \textbf{Internal Server Error} instead of a timeout -- the \texttt{connect()} call was blocked by the kernel before the connection could be established.

\begin{tcolorbox}[colback=green!5!white,colframe=green!75!black,title=Question]
\begin{itemize}
    \item Check the Falco logs after attempting the blocked attack. Does Falco generate any alert? Why or why not?
    \item What is the trade-off of blocking \texttt{connect()} globally? Can you think of a legitimate use case in \texttt{zipapp} that would require an outbound TCP connection?
    \item Restore the original deployment (without Seccomp) before proceeding:
\end{itemize}
\begin{lstlisting}[language=bash]
kubectl apply -f ../manifests/zipapp-deployment.yaml
kubectl rollout status deployment/zipapp
\end{lstlisting}
\end{tcolorbox}

\section{Blocking Ptrace}

Similarly, the ptrace attack can be blocked at the kernel level. Copy the second profile to the worker node, overwriting the previous \texttt{seccomp-profile.json} file so the deployment manifest can be applied without modifications:

\begin{lstlisting}[language=bash]
scp seccomp-profile-ptrace.json user@<worker-node-ip>:/tmp/
ssh user@<worker-node-ip> sudo mv /tmp/seccomp-profile-ptrace.json /var/lib/kubelet/seccomp/seccomp-profile.json
\end{lstlisting}

The profile \texttt{seccomp-profile.json} (copied from \texttt{seccomp-profile-ptrace.json}) blocks the \texttt{ptrace} syscall in the same manner. Apply the manifest to trigger a rollout and wait for it to complete:

\begin{lstlisting}[language=bash]
kubectl apply -f ../manifests/zipapp-deployment-seccomp.yaml
kubectl rollout status deployment/zipapp
\end{lstlisting}

Upload \texttt{ptrace\_payload.zip} and trigger the attack. You will again receive an Internal Server Error -- the \texttt{ptrace()} syscall returns \texttt{EPERM} and the binary exits immediately with an error.

\begin{tcolorbox}[colback=green!5!white,colframe=green!75!black,title=Question]
\begin{itemize}
    \item With Seccomp blocking \texttt{ptrace}, does Falco still detect the \textbf{attempt}? Check the logs and explain your finding.
    \item Reflect on the complementary roles of Falco and Seccomp. If you had to choose only one for a production environment, which would you prioritize and why?
\end{itemize}
\end{tcolorbox}

\section{Clean Up}

Once you have finished the lab, clean up all resources:

\begin{lstlisting}[language=bash]
# Restore the original deployment
kubectl apply -f ../manifests/zipapp-deployment.yaml

# Remove the Falco handler
kubectl delete deployment falco-handler
kubectl delete service falco-handler
kubectl delete clusterrolebinding falco-handler-deployment-binding
kubectl delete clusterrole falco-deployment-manager
kubectl delete serviceaccount falco-handler

# Remove the application and Falco
kubectl delete deployment zipapp
kubectl delete service zipapp
kubectl delete pod attacker --ignore-not-found
helm uninstall falco -n falco
kubectl delete namespace falco
\end{lstlisting}

